\documentclass{article}
\usepackage[T1]{fontenc}
\usepackage{lmodern}
\usepackage{graphicx}
\usepackage{amsmath,amssymb}
\usepackage[a4paper,margin=2cm]{geometry}
\usepackage[absolute,overlay]{textpos}
\setlength{\TPHorizModule}{1mm}
\setlength{\TPVertModule}{1mm}
\usepackage[indonesian]{babel}
\usepackage{hyperref}
\setlength{\emergencystretch}{2em}
% unit: Halaman 67

\begin{document}

% === Halaman 67 (llm) ===

Proses pembangkitan bit-bit acak di dalam A5/1 berdasarkan kunci sesi $K$ adalah sebagai berikut:

\begin{enumerate}
    \item Ketiga register pada awalnya diinisialisasi seluruh bitnya dengan 0. Selanjutnya dilakukan 64 putaran pertama, yaitu 64 bit kunci sesi $K$ dicampur dengan bit register berdasarkan skema berikut:
    
    Pada putaran $0 \le i < 64$,
    \begin{itemize}
        \item bit $K[i]$ ditambahkan ke bit $LSB$ dari setiap register $R$ dengan menggunakan operasi XOR:
        \[ R[0] = R[0] \oplus K[i] \]
        \item tiap register kemudian didetak
    \end{itemize}
\end{enumerate}

\end{document}
